% =============================================================================
%  The Grace Protocol: Dual-Track Input Validation for Graceful Cyber-Security
%
%  Author: Rafael D. De Paz
%  Date:   March 2026
% =============================================================================
\documentclass[12pt, a4paper]{article}

% --- Packages ---
\usepackage[utf8]{inputenc}
\usepackage[T1]{fontenc}
\usepackage{amsmath, amssymb, amsthm, mathtools}
\usepackage{geometry}
\usepackage{hyperref}
\usepackage{cleveref}
\usepackage{enumitem}
\usepackage{graphicx}
\usepackage{booktabs}
\usepackage{microtype}
\usepackage{natbib}

% --- Page Geometry ---
\geometry{margin=1in}

% --- Hyperref Setup ---
\hypersetup{
  colorlinks=true,
  linkcolor=blue,
  citecolor=blue,
  urlcolor=blue,
  pdftitle={Formal Paper},
  pdfauthor={Rafael D. De Paz}
}

% =============================================================================
\begin{document}

% --- Title ---
\title{
  \textbf{The Grace Protocol: Dual-Track Input Validation for Graceful Cyber-Security}
}
\author{
  Rafael D. De Paz \\
  \textit{Sovereign Architect} \\
  \texttt{me@rdepaz.com}
}
\date{March 2026}
\maketitle

% =============================================================================
% ABSTRACT
% =============================================================================
\begin{abstract}
In formal Language-Theoretic Security (LangSec), deterministic parsers enforce absolute security but often suffer from false positives on edge-cases. This paper introduces the Grace Protocol, a dual-track validation architecture. Track A (The Justiciary) applies strict, $O(1)$ determinism to reject malformed packets. Track B (The Intercessor) utilizes heuristic, Bayesian models to analyze non-malicious drift. We prove that routing edge cases through heuristic mercy rather than immediate terminal rejection creates a computationally superior Intrusion Detection System.
\end{abstract}

% --- Theorem Environments ---
\theoremstyle{plain}
\newtheorem{theorem}{Theorem}[section]
\newtheorem{proposition}[theorem]{Proposition}
\newtheorem{lemma}[theorem]{Lemma}

\section{Introduction}
These foundational principles operate on established invariants \cite{sassaman2011language,bratus2014exploits}.
Within the constraints of Language-Theoretic Security (LangSec), bounding the execution of network firewalls relies on parsing packet grammars. While strict, deterministic parsing prevents input-layer exploits (Zero-Day vulnerabilities), it introduces inherent brittleness, resulting in elevated false-positive packet rejection during semantic scaling or client state shifts.

This paper establishes the \textit{Grace Protocol}, a Dual-Track Computational Architecture that structurally divides the ingestion pipeline into a Justiciary Track (strict deterministic execution) and an Intercession Track (heuristic flexibility), proving that computational justice coupled with heuristic mercy yields optimal structural stability.

\section{Systemic Architecture: The Dual-Track Matrix}

Let the total input space of a network be defined as $\mathcal{I}$. 
Let the Acceptable Grammar defining secure, state-synchronizing packets be $G$. 

\subsection{Track A: The Justiciary (Deterministic Parser)}
The Justiciary acts as the primary firewall, operating over a strict Parsing Expression Grammar (PEG). It dictates whether an input string $s \in \mathcal{I}$ belongs to language $L(G)$.

\begin{theorem}[Justiciary Validation Speed]
By mapping $G$ strictly to a deterministic pushdown automaton, the validation time of $s$ becomes bounded by $\mathcal{O}(|s|)$, achieving an $O(1)$ constant overhead limit independent of network load.
\end{theorem}
The Justiciary enforces absolute justice. If $s \notin L(G)$, it is instantly marked for rejection. Let the set of rejected packets be $\mathcal{R}$.

\subsection{Track B: The Intercessor (Heuristic Mercy)}
Instead of terminating the connection for elements in $\mathcal{R}$, they are rerouted to the Intercessor (Track B). The Intercessor evaluates the \textit{Intent Boundary} of the packet using Bayesian heuristic networks.

Let $\mathcal{H}(s)$ be the probability that $s$ originated from a benign semantic shift rather than a malicious intent vector. 

\begin{proposition}[Optimization of Mercy]
For any rejected packet $s \in \mathcal{R}$, there exists an intervention threshold $\tau_h$ such that if $\mathcal{H}(s) > \tau_h$, dropping the packet produces more entropic shock to the network macro-node than dynamically adjusting the client state.
\end{proposition}

By evaluating:
\begin{equation}
  \Delta S_{\text{Grace}} = \mathcal{H}(s) \cdot S_{\text{recovery}} - (1 - \mathcal{H}(s)) \cdot S_{\text{breach}}
\end{equation}
The Grace Protocol acts entirely out of mathematical optimization. If $\Delta S_{\text{Grace}}$ yields a net thermodynamic negative (i.e., less entropy to forgive than to punish), the Intercessor applies dynamic translation, bridging the malformed packet into a safe canonical form.

\section{Conclusion}
Security architecture must mimic cosmological structure. A pure Justiciary framework achieves maximal security at the cost of structural sociability. A pure Intercession framework achieves scale at the cost of systemic collapse. The Grace Protocol balances these constraints, mathematically proving that computational forgiveness optimizes topological stability.

% =============================================================================
% REFERENCES
% =============================================================================

\vspace{2em}
\noindent\hrulefill
\vspace{1em}

\noindent\textbf{Cryptographic Lineage \& Validation}\\
This mathematical substrate has been cryptographically sealed and tracked on the global Sovereign Master Ledger to prevent retroactive editing and to verify the source authorship of Rafael D. De Paz. The immutable SHA-256 integrity checksums, formal PDF renderings, and lineage authorities can be verified at \url{https://rdepaz.com/research}.

\bibliographystyle{plainnat}
\bibliography{references}

\end{document}
